The Joukowsky ellipse family is the organizing complex-variable object in this section. Its exterior scale gives the logarithmic capacity of the confocal boundary, and that same scale provides a common logarithmic coordinate for multiplicative prime data, Lorentz composition, and geometric size invariants. Thus the finite register statements below are used only through the capacity, area, condition-number, and normalized-volume characters attached to Joukowsky ellipses. The hyperbolic identities below are used as a dictionary; the structural point is that Lorentz composition together with exact condition-number recovery forces this dictionary uniquely, up to reversing the rapidity orientation.

\begin{theorem}[G\"odel--Lorentz algebraization: Minkowski invariants and composition laws]\label{thm:group-jg-godel-lorentz-algebraization}
For \(L\ge 1\), define
\[
E_J(L):=\bigl\{\,L^{1/2}z+L^{-1/2}z^{-1}:|z|=1\,\bigr\}\subset\CC,
\]
\[
a(L):=L^{1/2}+L^{-1/2},
\qquad
b(L):=L^{1/2}-L^{-1/2}.
\]
For a prime register \(r\in\mathsf P\), write
\[
L:=N(r),
\qquad
E_J(r):=E_J(L),
\qquad
a(r):=a(L),
\qquad
b(r):=b(L),
\qquad
\Delta(r):=\frac12\log N(r).
\]
Define
\[
\Lambda(r):=
\begin{pmatrix}
\dfrac{a(r)}{2} & \dfrac{b(r)}{2}\\[2pt]
\dfrac{b(r)}{2} & \dfrac{a(r)}{2}
\end{pmatrix}\in\GL_2(\RR).
\]
Then:
\begin{enumerate}
  \item One has
  \[
  a(r)^2-b(r)^2=4,
  \qquad
  \det\Lambda(r)=1.
  \]
  Hence, whenever \(N(r)>1\), the curve \(E_J(r)\) is a confocal ellipse with focal distance \(2\); for \(N(r)=1\) it degenerates to the segment \([-2,2]\).

  \item Writing \(\Delta=\Delta(r)\), one has
  \[
  a(r)=2\cosh\Delta,
  \qquad
  b(r)=2\sinh\Delta,
  \]
  so \(\Lambda(r)\in \mathrm{SO}^+(1,1)\subset \mathrm{SL}_2(\RR)\), and the eigenvalues of \(\Lambda(r)\) are \(e^{\pm\Delta}\).

  \item The map \(\Lambda:\mathsf P\to \mathrm{SO}^+(1,1)\) is a commutative monoid homomorphism:
  \[
  \Lambda(r_1+r_2)=\Lambda(r_1)\Lambda(r_2).
  \]
  Equivalently, if \(a_i=a(r_i)\) and \(b_i=b(r_i)\), then
  \[
  a(r_1+r_2)=\frac{a_1a_2+b_1b_2}{2},
  \qquad
  b(r_1+r_2)=\frac{a_1b_2+b_1a_2}{2}.
  \]

  \item The spectral radius satisfies
  \[
  \log\rho\bigl(\Lambda(r)\bigr)=\Delta(r)=\frac12\log N(r).
  \]
\end{enumerate}
\end{theorem}

\begin{proof}
By direct calculation,
\[
a(r)^2-b(r)^2
=\bigl(L^{1/2}+L^{-1/2}\bigr)^2-\bigl(L^{1/2}-L^{-1/2}\bigr)^2
=4.
\]
Therefore \(\det\Lambda(r)=1\). If \(L>1\), the semiaxes of \(E_J(L)\) are \(a(L)\) and \(b(L)\), so the focal distance is \(\sqrt{a(L)^2-b(L)^2}=2\).

Writing \(L^{1/2}=e^\Delta\) gives
\[
a(r)=e^\Delta+e^{-\Delta}=2\cosh\Delta,
\qquad
b(r)=e^\Delta-e^{-\Delta}=2\sinh\Delta.
\]
Hence
\[
\Lambda(r)=
\begin{pmatrix}
\cosh\Delta & \sinh\Delta\\
\sinh\Delta & \cosh\Delta
\end{pmatrix},
\]
which belongs to \(\mathrm{SO}^+(1,1)\) and has eigenvalues \(\cosh\Delta\pm\sinh\Delta=e^{\pm\Delta}\).

Since \(N(r_1+r_2)=N(r_1)N(r_2)\), one has \(\Delta(r_1+r_2)=\Delta(r_1)+\Delta(r_2)\). The matrix identity in part (3) is therefore the standard addition law for hyperbolic functions, and the bilinear formulas for \((a,b)\) follow by reading off matrix entries. Finally, the spectral radius of \(\Lambda(r)\) is its larger eigenvalue \(e^{\Delta(r)}\), which proves part (4).
\end{proof}

\begin{theorem}[Rigidity of the G\"odel--Lorentz dictionary]\label{thm:group-jg-godel-lorentz-dictionary-rigidity}
For \(t\in\RR\), write
\[
B(t):=
\begin{pmatrix}
\cosh t&\sinh t\\
\sinh t&\cosh t
\end{pmatrix}
\in\mathrm{SO}^+(1,1).
\]
Let \(\Psi:\mathsf P\to\mathrm{SO}^+(1,1)\) be a monoid homomorphism. Assume that the condition number exactly recovers the G\"odel norm,
\[
\kappa_2(\Psi(r))=N(r)\qquad(r\in\mathsf P).
\]
Then either
\[
\Psi(r)=B\!\left(\frac12\log N(r)\right)\qquad(r\in\mathsf P)
\]
or
\[
\Psi(r)=B\!\left(-\frac12\log N(r)\right)\qquad(r\in\mathsf P).
\]
In particular, if \(\Psi(e_p)_{12}\ge 0\) for each prime generator \(e_p\), then \(\Psi=\Lambda\). Thus the assignment in Theorem \ref{thm:group-jg-godel-lorentz-algebraization} is not merely a reparameterization: it is the unique positive Lorentz representation whose singular-value distortion is exactly the multiplicative G\"odel code.
\end{theorem}

\begin{proof}
Every element of \(\mathrm{SO}^+(1,1)\) is uniquely of the form \(B(t)\) with \(t\in\RR\), and \(B(s)B(t)=B(s+t)\). Hence there is an additive map \(\tau:\mathsf P\to\RR\) such that \(\Psi(r)=B(\tau(r))\). Since the singular values of \(B(t)\) are \(e^t\) and \(e^{-t}\), its \(2\)-norm condition number is \(e^{2|t|}\). Applying the assumed condition-number identity to the prime generator \(e_p\) gives
\[
e^{2|\tau(e_p)|}=p,
\qquad\text{so}\qquad
\tau(e_p)=\varepsilon_p\frac12\log p
\]
with \(\varepsilon_p\in\{1,-1\}\).

It remains to show that the signs are all the same. If two prime generators \(e_p,e_q\) had opposite signs, then additivity would give
\[
2|\tau(e_p+e_q)|=|\varepsilon_p\log p+\varepsilon_q\log q|<\log p+
\log q=\log N(e_p+e_q),
\]
contradicting \(\kappa_2(\Psi(e_p+e_q))=N(e_p+e_q)\). Thus \(\varepsilon_p\) is independent of \(p\). Additivity now gives
\[
\tau(r)=\pm\frac12\sum_p r(p)\log p=\pm\frac12\log N(r),
\]
which proves the dichotomy. The positive-cone condition \(\Psi(e_p)_{12}\ge0\) selects the plus sign, because the off-diagonal entry of \(B(t)\) is \(\sinh t\).
\end{proof}

\begin{corollary}[Condition number identity]\label{cor:group-jg-godel-condition-number-identity}
In the notation of Theorem \ref{thm:group-jg-godel-lorentz-algebraization},
\[
\kappa_2\bigl(\Lambda(r)\bigr)=N(r).
\]
Equivalently,
\[
\log\kappa_2\bigl(\Lambda(r)\bigr)=\log N(r)=2\Delta(r).
\]
\end{corollary}

\begin{proof}
The matrix \(\Lambda(r)\) is symmetric with eigenvalues \(e^{\pm\Delta(r)}\). Its singular values are therefore \(e^{\pm\Delta(r)}\), and the \(2\)-norm condition number is
\[
\kappa_2\bigl(\Lambda(r)\bigr)=\frac{e^{\Delta(r)}}{e^{-\Delta(r)}}=e^{2\Delta(r)}=N(r).
\qedhere
\]
\end{proof}

\begin{theorem}[Area-preserving normalization and Einstein addition]\label{thm:group-jg-joukowsky-area-preserving-cayley}
For \(a>0\), let
\[
E(a):=\{(x,y)\in\RR^2:\ x^2/a^2+y^2a^2=1\}.
\]
If \(r\in\mathsf P\) satisfies \(N(r)>1\), define the area-preserving normalization
\[
\widehat E_J(r):=(a(r)b(r))^{-1/2}\cdot E_J(r).
\]
Then:
\begin{enumerate}
  \item \(\widehat E_J(r)\) belongs to the family \(E(a)\), with parameter
  \[
  \alpha(r):=\sqrt{\frac{a(r)}{b(r)}}
  =\sqrt{\frac{L+1}{L-1}}
  =\sqrt{\coth(\Delta(r))}\in(1,\infty).
  \]

  \item If \(u(r):=\alpha(r)^2=\coth(\Delta(r))\), then
  \[
  u(r_1+r_2)=u(r_1)\odot u(r_2),
  \qquad
  u\odot v:=\frac{uv+1}{u+v}.
  \]
  Equivalently, if \(t(r):=u(r)^{-1}=\tanh(\Delta(r))\in(0,1)\), then
  \[
  t(r_1+r_2)=\frac{t(r_1)+t(r_2)}{1+t(r_1)t(r_2)}.
  \]
\end{enumerate}
\end{theorem}

\begin{proof}
The normalized semiaxes of \(\widehat E_J(r)\) are
\[
\hat a=\frac{a(r)}{\sqrt{a(r)b(r)}}=\sqrt{\frac{a(r)}{b(r)}},
\qquad
\hat b=\frac{b(r)}{\sqrt{a(r)b(r)}}=\sqrt{\frac{b(r)}{a(r)}}=\hat a^{-1}.
\]
Hence \(\widehat E_J(r)\) is described by \(x^2/\hat a^2+y^2\hat a^2=1\), so it belongs to the area-preserving family \(E(a)\) with parameter \(\alpha(r)=\hat a\).

Now
\[
\frac{a(r)}{b(r)}
=\frac{L^{1/2}+L^{-1/2}}{L^{1/2}-L^{-1/2}}
=\frac{L+1}{L-1},
\]
and since \(L=e^{2\Delta(r)}\), this equals \(\coth(\Delta(r))\). The composition laws for \(u\) and \(t\) are the standard identities
\[
\coth(\Delta_1+\Delta_2)
=\frac{\coth\Delta_1\,\coth\Delta_2+1}{\coth\Delta_1+\coth\Delta_2},
\]
\[
\tanh(\Delta_1+\Delta_2)
=\frac{\tanh\Delta_1+\tanh\Delta_2}{1+\tanh\Delta_1\tanh\Delta_2},
\]
together with \(\Delta(r_1+r_2)=\Delta(r_1)+\Delta(r_2)\).
\end{proof}

\begin{proposition}[Primorial lower bounds for \(p\)-ary G\"odel encodings]\label{prop:group-jg-boundary-primorial-ellipse-bound}
Let \(p\) be a prime and \(b\ge 1\). Choose distinct primes \(q_1,\dots,q_b\), and define the scalar G\"odel encoding
\[
g(\varepsilon_1,\dots,\varepsilon_b):=\prod_{i=1}^{b}q_i^{\varepsilon_i}
\qquad
(\varepsilon_i\in\{0,1,\dots,p-1\}).
\]
Let
\[
L_{\max}:=\max_{\varepsilon}g(\varepsilon)=\prod_{i=1}^{b}q_i^{p-1},
\]
and write \(p_b\#:=\prod_{i=1}^{b}p_i\) for the \(b\)th primorial. Then:
\begin{enumerate}
  \item
  \[
  L_{\max}\ge (p_b\#)^{p-1},
  \]
  with equality if and only if \(\{q_1,\dots,q_b\}=\{p_1,\dots,p_b\}\).

  \item The associated rapidity
  \[
  \Delta_{\max}:=\frac12\log L_{\max}
  \]
  satisfies
  \[
  \Delta_{\max}\ge \frac{p-1}{2}\log(p_b\#).
  \]

  \item Let
  \[
  a_{\max}:=a(L_{\max}),
  \qquad
  b_{\max}:=b(L_{\max}).
  \]
  Then
  \[
  a_{\max}\ge (p_b\#)^{\frac{p-1}{2}}+(p_b\#)^{\frac{1-p}{2}},
  \qquad
  b_{\max}\ge (p_b\#)^{\frac{p-1}{2}}-(p_b\#)^{\frac{1-p}{2}},
  \]
  and
  \[
  \mathrm{Area}\bigl(E_J(L_{\max})\bigr)
  =\pi\bigl(L_{\max}-L_{\max}^{-1}\bigr)
  \ge
  \pi\bigl((p_b\#)^{p-1}-(p_b\#)^{1-p}\bigr).
  \]
  The eccentricity \(e_{\max}:=2/a_{\max}\) therefore satisfies
  \[
  e_{\max}\le
  \frac{2\,(p_b\#)^{(p-1)/2}}{(p_b\#)^{p-1}+1}.
  \]

  \item The normalized parameter
  \[
  \alpha_{\max}:=\sqrt{\frac{L_{\max}+1}{L_{\max}-1}}
  \]
  satisfies
  \[
  \alpha_{\max}\le
  \sqrt{\frac{(p_b\#)^{p-1}+1}{(p_b\#)^{p-1}-1}},
  \]
  so this upper bound decreases monotonically to \(1\) as \(b\to\infty\).
\end{enumerate}
\end{proposition}

\begin{proof}
After reordering, write \(q_{(1)}<\cdots<q_{(b)}\). Since the \(i\)th smallest prime is \(p_i\), one has \(q_{(i)}\ge p_i\) for each \(i\). Therefore
\[
L_{\max}
=\prod_{i=1}^{b}q_{(i)}^{p-1}
\ge
\prod_{i=1}^{b}p_i^{p-1}
=(p_b\#)^{p-1},
\]
with equality if and only if \(q_{(i)}=p_i\) for all \(i\). The bound for \(\Delta_{\max}\) follows by taking logarithms.

The remaining statements follow from the monotonicity, for \(L>1\), of the functions
\[
a(L)=L^{1/2}+L^{-1/2},
\qquad
b(L)=L^{1/2}-L^{-1/2},
\]
\[
\mathrm{Area}\bigl(E_J(L)\bigr)=\pi a(L)b(L)=\pi(L-L^{-1}),
\qquad
e(L)=\frac{2}{a(L)},
\qquad
\alpha(L)^2=\frac{L+1}{L-1}.
\]
Indeed, \(a\), \(b\), and the area are increasing in \(L\), while \(e\) and \(\alpha\) are decreasing in \(L\). Applying these monotonicity facts to the lower bound \(L_{\max}\ge (p_b\#)^{p-1}\) gives all displayed inequalities.
\end{proof}

\begin{corollary}[Primorial forcing for Lorentz condition numbers]\label{cor:group-jg-dyadic-boundary-godel-ellipse-illconditioning}
Under the hypotheses of Proposition \ref{prop:group-jg-boundary-primorial-ellipse-bound}, let \(r_{\max}\in\mathsf P\) be any prime register satisfying \(N(r_{\max})=L_{\max}\). Then
\[
\kappa_2\bigl(\Lambda(r_{\max})\bigr)=L_{\max}\ge (p_b\#)^{p-1},
\]
and therefore
\[
\log\kappa_2\bigl(\Lambda(r_{\max})\bigr)\ge (p-1)\log(p_b\#).
\]
Equality holds if and only if the primes used in the encoding are precisely the first \(b\) primes.
\end{corollary}

\begin{proof}
The condition-number identity gives
\[
\kappa_2\bigl(\Lambda(r_{\max})\bigr)=N(r_{\max})=L_{\max}.
\]
The claim now follows immediately from Proposition \ref{prop:group-jg-boundary-primorial-ellipse-bound}(1).
\end{proof}

\begin{corollary}[Capacity, modulus, and rapidity]\label{cor:group-jg-capacity-modulus-koenigs}
Let \(r\in\mathsf P\) satisfy \(N(r)>1\). Then the logarithmic capacity of the Joukowsky ellipse \(E_J(r)\) is
\[
\mathrm{cap}\bigl(E_J(r)\bigr)=N(r)^{1/2},
\qquad
\log N(r)=2\log\mathrm{cap}\bigl(E_J(r)\bigr).
\]
Moreover, for the annulus
\[
A_{N(r)^{1/2}}:=\{z\in\CC:\ 1<|z|<N(r)^{1/2}\},
\]
its conformal modulus is
\[
\mathrm{mod}\bigl(A_{N(r)^{1/2}}\bigr)
=\frac{1}{2\pi}\Delta(r),
\qquad
\log N(r)=4\pi\,\mathrm{mod}\bigl(A_{N(r)^{1/2}}\bigr).
\]
\end{corollary}

\begin{proof}
The annulus formula
\[
\mathrm{mod}\{\,\rho<|z|<R\,\}=\frac{1}{2\pi}\log(R/\rho)
\]
gives the second statement with \(\rho=1\) and \(R=N(r)^{1/2}\). For the capacity, the exterior Joukowsky map
\[
w=L^{1/2}z+L^{-1/2}z^{-1}
\]
has expansion \(w=L^{1/2}z+O(z^{-1})\) at infinity, so the conformal radius of the exterior domain is \(L^{1/2}=N(r)^{1/2}\). This is exactly the logarithmic capacity of \(E_J(r)\).
\end{proof}

The next result is used as the volume interface for the primorial codebook theorem, not as an independent excursion into higher-dimensional convex geometry.  A one-dimensional Joukowsky ellipse records only the scalar norm \(N(r)\).  The diagonal ellipsoid below keeps the prime coordinates separated until the final determinant is taken; this is exactly what lets the pairwise-coprime codebook constraint turn into a product-volume minimization in Theorem \ref{prop:group-jg-min-total-godel-cost-primorial}.

\begin{theorem}[A higher-dimensional ellipsoid functor for prime registers]\label{thm:group-jg-godel-ellipsoid-volume}
Let \(P_0=\{p_1,\dots,p_d\}\) be a finite set of primes and let \(r\in\mathsf P\) be supported in \(P_0\). Write \(r_i:=r(p_i)\). Define the diagonal ellipsoid
\[
\mathcal E_d(r):=\left\{x\in\RR^d:\ \sum_{i=1}^{d}\frac{x_i^2}{p_i^{r_i}}\le 1\right\}.
\]
Let
\[
A_r:=\operatorname{diag}\bigl(p_1^{r_1/2},\dots,p_d^{r_d/2}\bigr).
\]
Then:
\begin{enumerate}
  \item \(\mathcal E_d(r)=A_rB_d\), and
\[
\mathrm{Vol}\bigl(\mathcal E_d(r)\bigr)
=\mathrm{Vol}(B_d)\prod_{i=1}^{d}p_i^{r_i/2}
=\mathrm{Vol}(B_d)\,\sqrt{N(r)}.
\]
Consequently,
\[
\log\mathrm{Vol}\bigl(\mathcal E_d(r)\bigr)-\log\mathrm{Vol}(B_d)
=\frac12\log N(r)=\Delta(r).
\]

  \item The assignment \(r\mapsto A_r\) is a faithful commutative monoid representation of registers supported in \(P_0\):
  \[
  A_{r+s}=A_rA_s,
  \qquad
  \det A_r=\sqrt{N(r)}.
  \]
  Hence the normalized ellipsoid-volume character
  \[
  V_d(r):=\frac{\mathrm{Vol}(\mathcal E_d(r))}{\mathrm{Vol}(B_d)}
  \]
  satisfies
  \[
  V_d(r+s)=V_d(r)V_d(s),
  \qquad
  \log V_d(r)=\Delta(r).
  \]

  \item If \(r_1,
\dots,r_n\) are supported in \(P_0\), then
  \[
  \prod_{k=1}^{n}\mathrm{Vol}\bigl(\mathcal E_d(r_k)\bigr)
  =\mathrm{Vol}(B_d)^n\exp\!\left(\frac12\sum_{k=1}^{n}\log N(r_k)\right).
  \]
  In particular, for a G\"odel-decomposable codebook \(\iota:X\hookrightarrow\NN_{\mathrm{sf},>1}\) with exponent registers \(r_x\), the ellipsoid-volume product is a strictly increasing function of the total G\"odel cost \(C(\iota)=\sum_x\log\iota(x)\). Thus the ellipsoid functor descends the prime-register decomposition to the primorial capacity minimization problem without losing the coordinatewise coprimality data before the determinant is evaluated.
\end{enumerate}
\end{theorem}

\begin{proof}
The ellipsoid \(\mathcal E_d(r)\) is the image of the unit ball \(B_d\) under the diagonal linear map
\[
A_r(x_1,\dots,x_d)=\bigl(p_1^{r_1/2}x_1,\dots,p_d^{r_d/2}x_d\bigr).
\]
Therefore
\[
\mathrm{Vol}\bigl(\mathcal E_d(r)\bigr)=|\det A_r|\,\mathrm{Vol}(B_d)
=\left(\prod_{i=1}^{d}p_i^{r_i/2}\right)\mathrm{Vol}(B_d)
=\sqrt{N(r)}\,\mathrm{Vol}(B_d).
\]
Taking logarithms yields the final identity.
The equality \(A_{r+s}=A_rA_s\) follows coordinate by coordinate from
\[
p_i^{(r_i+s_i)/2}=p_i^{r_i/2}p_i^{s_i/2},
\]
and \(\det A_r=\prod_i p_i^{r_i/2}=\sqrt{N(r)}\). If \(A_r=A_s\), then every diagonal entry agrees, so \(p_i^{r_i/2}=p_i^{s_i/2}\) and hence \(r_i=s_i\) for all \(i\); this proves faithfulness on registers supported in \(P_0\). Dividing the volume formula by \(\mathrm{Vol}(B_d)\) gives \(V_d(r)=\det A_r\), so multiplicativity and \(\log V_d(r)=\Delta(r)\) follow.

For registers \(r_1,\dots,r_n\), multiplying the volume formula gives
\[
\prod_{k=1}^{n}\mathrm{Vol}\bigl(\mathcal E_d(r_k)\bigr)
=\mathrm{Vol}(B_d)^n\prod_{k=1}^{n}\sqrt{N(r_k)}
=\mathrm{Vol}(B_d)^n\exp\!\left(\frac12\sum_{k=1}^{n}\log N(r_k)\right).
\]
If the registers come from a G\"odel-decomposable codebook, then \(N(r_x)=\iota(x)\), so the exponent is \(C(\iota)/2\). Since the exponential function is strictly increasing, minimizing the ellipsoid-volume product is equivalent to minimizing the total G\"odel cost. This is the precise interface used in Theorem \ref{prop:group-jg-min-total-godel-cost-primorial}.
\end{proof}

\begin{corollary}[Volume gap and ellipse area]\label{cor:group-jg-ellipse-area-ellipsoid-gap}
Let \(r\in\mathsf P\) satisfy \(N(r)>1\). Then
\[
\mathrm{Area}\bigl(E_J(r)\bigr)
=\pi\bigl(e^{2\Delta(r)}-e^{-2\Delta(r)}\bigr)
=\pi\bigl(N(r)-N(r)^{-1}\bigr).
\]
Hence
\[
\log\mathrm{Area}\bigl(E_J(r)\bigr)
=2\Delta(r)+\log\pi+\log\bigl(1-e^{-4\Delta(r)}\bigr),
\]
and in particular
\[
\log\mathrm{Area}\bigl(E_J(r)\bigr)=2\Delta(r)+\log\pi+o(1)
\qquad (\Delta(r)\to\infty).
\]
\end{corollary}

\begin{proof}
By Theorem \ref{thm:group-jg-godel-lorentz-algebraization}, the semiaxes of \(E_J(r)\) are
\[
a(r)=e^{\Delta(r)}+e^{-\Delta(r)},
\qquad
b(r)=e^{\Delta(r)}-e^{-\Delta(r)}.
\]
Therefore
\[
\mathrm{Area}\bigl(E_J(r)\bigr)
=\pi a(r)b(r)
=\pi\bigl(e^{2\Delta(r)}-e^{-2\Delta(r)}\bigr)
=\pi\bigl(N(r)-N(r)^{-1}\bigr).
\]
The logarithmic formula follows by factoring out \(e^{2\Delta(r)}\), and the asymptotic statement is immediate because \(\log(1-e^{-4\Delta(r)})\to 0\).
\end{proof}


\begin{theorem}[Capacity-normalized primorial extremality]\label{thm:group-jg-capacity-normalized-primorial-extremality}
Let \(X\) be a finite set with \(|X|=n\ge1\), and let
\[
\iota:X\hookrightarrow \NN_{\mathrm{sf},>1}
\]
be G\"odel-decomposable.  For each \(x\in X\), let \(r_x\in\mathsf P\) be the exponent register of \(\iota(x)\), so that \(N(r_x)=\iota(x)\). Define the normalized ellipsoid-volume character
\[
\mathcal V(\iota):=\prod_{x\in X}\frac{\mathrm{Vol}(\mathcal E_{D_x}(r_x))}{\mathrm{Vol}(B_{D_x})},
\qquad D_x:=|\operatorname{supp}(r_x)|,
\]
where \(\mathcal E_{D_x}(r_x)\) is formed in the coordinate space indexed by \(\operatorname{supp}(r_x)\). Define also the Joukowsky capacity product
\[
\mathcal C_J(\iota):=\prod_{x\in X}\mathrm{cap}\bigl(E_J(r_x)\bigr).
\]
Then
\[
\mathcal V(\iota)=\mathcal C_J(\iota)=\sqrt{\prod_{x\in X}\iota(x)},
\]
and therefore
\[
\mathcal V(\iota)\ge\sqrt{p_n\#},
\qquad
\mathcal C_J(\iota)\ge\sqrt{p_n\#}.
\]
Equality in either inequality holds if and only if
\[
\iota(X)=\{p_1,\dots,p_n\}.
\]
Moreover, if a finite ambient prime set \(P_0\) with \(|P_0|=D\) is fixed in advance and every competitor is required to have support in \(P_0\), then the unnormalized fixed-dimensional product satisfies
\[
\prod_{x\in X}\mathrm{Vol}\bigl(\mathcal E_D(r_x)\bigr)
=\mathrm{Vol}(B_D)^n\,\mathcal V(\iota),
\]
so the same primorial codebook is the unique minimizer whenever \(\{p_1,\dots,p_n\}\subset P_0\); if \(P_0\) contains at least \(n\) primes but omits one of the first \(n\) primes, the minimizer is instead the set of the \(n\) smallest primes contained in \(P_0\).
\end{theorem}

\begin{proof}
The volume formula in Theorem \ref{thm:group-jg-godel-ellipsoid-volume}, applied in the coordinate space of \(\operatorname{supp}(r_x)\), gives
\[
\frac{\mathrm{Vol}(\mathcal E_{D_x}(r_x))}{\mathrm{Vol}(B_{D_x})}
=\sqrt{N(r_x)}=\sqrt{\iota(x)}.
\]
Corollary \ref{cor:group-jg-capacity-modulus-koenigs} gives
\[
\mathrm{cap}\bigl(E_J(r_x)\bigr)=\sqrt{N(r_x)}=\sqrt{\iota(x)}.
\]
Multiplying over \(x\in X\) proves
\[
\mathcal V(\iota)=\mathcal C_J(\iota)=\sqrt{\prod_x\iota(x)}.
\]

It remains to minimize \(\prod_x\iota(x)\). For each \(x\), choose a prime divisor \(q_x\mid\iota(x)\). Pairwise coprimality of the image of \(\iota\) makes the primes \(q_x\) distinct. Ordering them as \(q_{(1)}<\cdots<q_{(n)}\), we have \(q_{(i)}\ge p_i\), hence
\[
\prod_{x\in X}\iota(x)\ge \prod_{x\in X}q_x\ge\prod_{i=1}^{n}p_i=p_n\#.
\]
Equality forces every \(\iota(x)\) to equal its chosen prime divisor and forces \(\{q_x:x\in X\}=\{p_1,\dots,p_n\}\), which is exactly the displayed primorial codebook. Conversely that codebook is admissible and attains equality.

For the fixed ambient statement, Theorem \ref{thm:group-jg-godel-ellipsoid-volume} gives
\[
\mathrm{Vol}\bigl(\mathcal E_D(r_x)\bigr)=\mathrm{Vol}(B_D)\sqrt{\iota(x)}
\]
with the same dimension \(D\) for every competitor, so the unnormalized product differs from \(\mathcal V(\iota)\) by the fixed factor \(\mathrm{Vol}(B_D)^n\). If \(P_0\) contains the first \(n\) primes, the preceding equality case applies. If \(P_0\) contains at least \(n\) primes but not the first \(n\) primes, the same ordering argument performed inside the ordered list of primes in \(P_0\) gives the product of the \(n\) smallest available primes, and equality again requires the codebook to consist exactly of those primes.
\end{proof}

\begin{theorem}[Primorial minimization for G\"odel-decomposable codebooks]\label{prop:group-jg-min-total-godel-cost-primorial}
Let \(X\) be a finite set with \(|X|=n\ge 1\), and let
\[
\iota:X\hookrightarrow \NN_{\mathrm{sf},>1}
\]
be a G\"odel-decomposable codebook in the sense of Definition \ref{def:godel-decomposable-codebook}. Let \(C(\iota)\) denote its total G\"odel cost. Then:
\begin{enumerate}
  \item \textbf{Minimal total cost.}
  Among all G\"odel-decomposable codebooks,
  \[
  \min_{\iota}C(\iota)=\log(p_n\#),
  \]
  and equality holds precisely when \(\iota(X)=\{p_1,\dots,p_n\}\).

  \item \textbf{Minimal normalized and fixed-ambient ellipsoid-volume product.}
  For the exponent registers \(r_x\) of \(\iota(x)\), the normalized character
  \[
  \prod_{x\in X}\frac{\mathrm{Vol}(\mathcal E_{D_x}(r_x))}{\mathrm{Vol}(B_{D_x})},
  \qquad D_x:=|\operatorname{supp}(r_x)|,
  \]
  is minimized precisely by the same primorial codebook. Equivalently, if one fixes in advance a finite ambient prime set \(P_0\) of cardinality \(D\) containing the first \(n\) primes and allows only codebooks supported in \(P_0\), then
  \[
  \prod_{x\in X}\mathrm{Vol}\bigl(\mathcal E_D(r_x)\bigr)
  =\mathrm{Vol}(B_D)^n\sqrt{\prod_{x\in X}\iota(x)}
  \]
  is minimized precisely by \(\iota(X)=\{p_1,\dots,p_n\}\).

  \item \textbf{Minimal ellipse-area product.}
  In the same notation,
  \[
  \prod_{x\in X}\mathrm{Area}\bigl(E_J(r_x)\bigr)
  =
  \pi^n\prod_{x\in X}\bigl(\iota(x)-\iota(x)^{-1}\bigr),
  \]
  and this product is again minimized precisely by \(\iota(X)=\{p_1,\dots,p_n\}\).
\end{enumerate}
\end{theorem}

\begin{proof}
\textbf{(1)} For each \(x\in X\), choose a prime divisor \(q_x\mid \iota(x)\). Because the image of \(\iota\) is pairwise coprime, the primes \(q_x\) are distinct. After ordering them increasingly as \(q_{(1)}<\cdots<q_{(n)}\), one has \(q_{(i)}\ge p_i\), so
\[
\prod_{x\in X}\iota(x)\ge \prod_{x\in X}q_x\ge \prod_{i=1}^{n}p_i=p_n\#.
\]
Taking logarithms gives \(C(\iota)\ge \log(p_n\#)\). Equality holds for \(\iota(X)=\{p_1,\dots,p_n\}\), and if equality occurs then every step above must be sharp, forcing \(\iota(x)=q_x\) and \(\{q_x\}=\{p_1,\dots,p_n\}\).

\textbf{(2)} This is the normalized and fixed-ambient specialization of Theorem \ref{thm:group-jg-capacity-normalized-primorial-extremality}. The fixed-dimensional factor \(\mathrm{Vol}(B_D)^n\) is constant only because the ambient prime set \(P_0\) is fixed before the minimization begins.

\textbf{(3)} By Corollary \ref{cor:group-jg-ellipse-area-ellipsoid-gap},
\[
\mathrm{Area}\bigl(E_J(r_x)\bigr)=\pi\bigl(\iota(x)-\iota(x)^{-1}\bigr).
\]
The function \(h(t):=t-t^{-1}\) is strictly increasing for \(t>1\), since \(h'(t)=1+t^{-2}>0\). Therefore
\[
\prod_{x\in X}\bigl(\iota(x)-\iota(x)^{-1}\bigr)
\ge
\prod_{x\in X}\bigl(q_x-q_x^{-1}\bigr)
\ge
\prod_{i=1}^{n}\bigl(p_i-p_i^{-1}\bigr),
\]
with equality only for the primorial codebook from part (1). This proves the minimality of the area product.
\end{proof}

\begin{remark}[Necessity of the coprimality constraint]\label{rem:group-jg-min-total-godel-cost-coprime-necessity}
The primorial minimization theorem is specific to the G\"odel-decomposable class. If one merely asks for a set of distinct square-free integers, shared prime factors can reduce the total product. For instance,
\[
\{2,3,5,6\}\subset \NN_{\mathrm{sf},>1}
\]
consists of distinct square-free integers, but its product \(180\) is smaller than \(p_4\#=210\). Pairwise coprimality is therefore the correct structural constraint for multiplicative decomposability.
\end{remark}

\endinput


